% GENERATED FILE. Edit canonical lessons, then run npm run generate:books.
% canonical-source-hash: fnv1a64:4afb2ddb7987ba77
% canonical-lessons: GU-C221-dukal, GU-C221-kharvu, GU-C221-ag, GU-C221-salagvu, GU-C221-bhukamp

\chapter{Drought fire and earthquake}
\label{ch:gu-drought-fire-and-earthquake}
\hlchaptermodality{221}

\begin{chapteropening}
\textbf{By the end of this chapter:} \emph{I can say what a hard year will bring.}

Complete the last lesson of chapter 221: I can say what a hard year will bring.
\end{chapteropening}

\section[\texorpdfstring{\gu{દુકાળ} (dukāḷ)}{દુકાળ (dukāḷ)}]{\gu{દુકાળ} (dukāḷ) — drought, famine}

\label{lesson:GU-C221-dukal}



\begin{quote}
Before the new one: say the Gujarati for to recede, then the Gujarati for to clear up.
\end{quote}

\subsection*{You'll want to know: \gu{દુકાળ}}
\textbf{\gu{દુકાળ}} — \emph{dukāḷ} — \textquotedblleft{}drought, famine\textquotedblright{}.

\begin{quote}
\textbf{\gu{આ} \gu{વરસે} \gu{દુકાળ} \gu{પડશે}.} — \emph{ā varase dukāḷ paḍśe} — \textquotedblleft{}There will be a drought this year.\textquotedblright{}
\end{quote}

A masculine noun.

\begin{cousinweb}
\textbf{\gu{દુકાળ}} (\emph{dukāḷ}) is a tadbhava. It comes from Sanskrit \emph{duṣkāla}, \textquotedblleft{}evil time\textquotedblright{}, built from \emph{duṣ-}, \textquotedblleft{}bad\textquotedblright{}, and \emph{kāla}, \textquotedblleft{}time\textquotedblright{}; CDIAL glosses the Gujarati word \textquotedblleft{}famine\textquotedblright{} (CDIAL 6470).
\end{cousinweb}

\begin{grammarlens}[title={What you've built}]
One more word for everyday life at A2.
\end{grammarlens}

\subsection*{Your turn}
\begin{itemize}
\raggedright
  \item \emph{Say it:} \emph{dukāḷ}
  \item \emph{Say it:} \emph{dukāḷ}, once more
  \item \emph{Say it:} \emph{ā varase dukāḷ paḍśe} — There will be a drought this year.
  \item \emph{Recall:} say the Gujarati for to recede, then the Gujarati for to clear up, then say \emph{dukāḷ} again
\end{itemize}

\begin{tcolorbox}[breakable,colback=teal!4,colframe=teal!35!black,title={Before you move on}]
What does \gu{દુકાળ} mean? (\textquotedblleft{}Drought, famine\textquotedblright{}.) Say it once more.
\end{tcolorbox}
\section[\texorpdfstring{\gu{ખરવું} (kharvũ)}{ખરવું (kharvũ)}]{\gu{ખરવું} (kharvũ) — to fall off, to shed}

\label{lesson:GU-C221-kharvu}



\begin{quote}
Before the new one: say the Gujarati for to clear up, then the Gujarati for drought.
\end{quote}

\subsection*{You'll want to know: \gu{ખરવું}}
\textbf{\gu{ખરવું}} — \emph{kharvũ} — \textquotedblleft{}to fall off, to shed\textquotedblright{}.

\begin{quote}
\textbf{\gu{પાનખરમાં} \gu{પાન} \gu{ખરશે}.} — \emph{pānkharmā̃ pān kharśe} — \textquotedblleft{}Leaves will fall in autumn.\textquotedblright{}
\end{quote}

\begin{cousinweb}
\textbf{\gu{ખરવું}} (\emph{kharvũ}) is a tadbhava. It comes from Sanskrit \emph{kṣarati}, \textquotedblleft{}trickles; melts away\textquotedblright{}, and CDIAL glosses the Gujarati verb \textquotedblleft{}to fall, fade\textquotedblright{} (CDIAL 3663).
\end{cousinweb}

\begin{grammarlens}[title={What you've built}]
One more word for everyday life at A2.
\end{grammarlens}

\subsection*{Your turn}
\begin{itemize}
\raggedright
  \item \emph{Say it:} \emph{kharvũ}
  \item \emph{Say it:} \emph{kharvũ}, once more
  \item \emph{Say it:} \emph{pānkharmā̃ pān kharśe} — Leaves will fall in autumn.
  \item \emph{Recall:} say the Gujarati for to clear up, then the Gujarati for drought, then say \emph{kharvũ} again
\end{itemize}

\begin{tcolorbox}[breakable,colback=teal!4,colframe=teal!35!black,title={Before you move on}]
What does \gu{ખરવું} mean? (\textquotedblleft{}To fall off, to shed\textquotedblright{}.) Say it once more.
\end{tcolorbox}
\section[\texorpdfstring{\gu{આગ} (āg)}{આગ (āg)}]{\gu{આગ} (āg) — fire}

\label{lesson:GU-C221-ag}



\begin{quote}
Before the new one: say the Gujarati for drought, then the Gujarati for to fall off.
\end{quote}

\subsection*{You'll want to know: \gu{આગ}}
\textbf{\gu{આગ}} — \emph{āg} — \textquotedblleft{}fire\textquotedblright{}.

\begin{quote}
\textbf{\gu{જંગલમાં} \gu{આગ} \gu{ફેલાશે}.} — \emph{jangalmā̃ āg phelāśe} — \textquotedblleft{}The fire will spread in the forest.\textquotedblright{}
\end{quote}

A feminine noun.

\begin{cousinweb}
\textbf{\gu{આગ}} (\emph{āg}) is a tadbhava, a word inherited from Sanskrit and reshaped by sound change. It comes from \emph{agni}, \textquotedblleft{}fire\textquotedblright{} (CDIAL 55).
\end{cousinweb}

\begin{grammarlens}[title={What you've built}]
One more word for everyday life at A2.
\end{grammarlens}

\subsection*{Your turn}
\begin{itemize}
\raggedright
  \item \emph{Say it:} \emph{āg}
  \item \emph{Say it:} \emph{āg}, once more
  \item \emph{Say it:} \emph{jangalmā̃ āg phelāśe} — The fire will spread in the forest.
  \item \emph{Recall:} say the Gujarati for drought, then the Gujarati for to fall off, then say \emph{āg} again
\end{itemize}

\begin{tcolorbox}[breakable,colback=teal!4,colframe=teal!35!black,title={Before you move on}]
What does \gu{આગ} mean? (\textquotedblleft{}Fire\textquotedblright{}.) Say it once more.
\end{tcolorbox}
\section[\texorpdfstring{\gu{સળગવું} (saḷagvũ)}{સળગવું (saḷagvũ)}]{\gu{સળગવું} (saḷagvũ) — to catch fire, to burn}

\label{lesson:GU-C221-salagvu}



\begin{quote}
Before the new one: say the Gujarati for to fall off, then the Gujarati for fire.
\end{quote}

\subsection*{You'll want to know: \gu{સળગવું}}
\textbf{\gu{સળગવું}} — \emph{saḷagvũ} — \textquotedblleft{}to catch fire, to burn\textquotedblright{}.

\begin{quote}
\textbf{\gu{સૂકું} \gu{ઘાસ} \gu{તરત} \gu{સળગશે}.} — \emph{sūkũ ghās tarat saḷagśe} — \textquotedblleft{}Dry grass will catch fire at once.\textquotedblright{}
\end{quote}

\begin{cousinweb}
\textbf{\gu{સળગવું}} (\emph{saḷagvũ}) is a tadbhava. CDIAL lists it, glossed \textquotedblleft{}to take fire\textquotedblright{}, under Sanskrit \emph{saṁlagna}, \textquotedblleft{}closely attached\textquotedblright{}, so the fire is what takes hold (CDIAL 12999).
\end{cousinweb}

\begin{grammarlens}[title={What you've built}]
One more word for everyday life at A2.
\end{grammarlens}

\subsection*{Your turn}
\begin{itemize}
\raggedright
  \item \emph{Say it:} \emph{saḷagvũ}
  \item \emph{Say it:} \emph{saḷagvũ}, once more
  \item \emph{Say it:} \emph{sūkũ ghās tarat saḷagśe} — Dry grass will catch fire at once.
  \item \emph{Recall:} say the Gujarati for to fall off, then the Gujarati for fire, then say \emph{saḷagvũ} again
\end{itemize}

\begin{tcolorbox}[breakable,colback=teal!4,colframe=teal!35!black,title={Before you move on}]
What does \gu{સળગવું} mean? (\textquotedblleft{}To catch fire, to burn\textquotedblright{}.) Say it once more.
\end{tcolorbox}
\section[\texorpdfstring{\gu{ભૂકંપ} (bhūkamp)}{ભૂકંપ (bhūkamp)}]{\gu{ભૂકંપ} (bhūkamp) — earthquake}

\label{lesson:GU-C221-bhukamp}



\begin{quote}
Before the new one: say the Gujarati for fire, then the Gujarati for to catch fire.
\end{quote}

\subsection*{You'll want to know: \gu{ભૂકંપ}}
\textbf{\gu{ભૂકંપ}} — \emph{bhūkamp} — \textquotedblleft{}earthquake\textquotedblright{}.

\begin{quote}
\textbf{\gu{કદાચ} \gu{અહીં} \gu{ભૂકંપ} \gu{આવશે}.} — \emph{kadāc ahī̃ bhūkamp āvśe} — \textquotedblleft{}Perhaps an earthquake will come here.\textquotedblright{}
\end{quote}

A masculine noun.

\begin{cousinweb}
\textbf{\gu{ભૂકંપ}} (\emph{bhūkamp}) is a learned loan, a tatsama, built from Sanskrit \emph{bhū}, \textquotedblleft{}earth\textquotedblright{} (CDIAL 9546), and \emph{kampa}, \textquotedblleft{}tremor\textquotedblright{} (CDIAL 2765), so the earth itself shakes.
\end{cousinweb}

\begin{grammarlens}[title={What you've built}]
One more word for everyday life at A2.
\end{grammarlens}

\subsection*{Your turn}
\begin{itemize}
\raggedright
  \item \emph{Say it:} \emph{bhūkamp}
  \item \emph{Say it:} \emph{bhūkamp}, once more
  \item \emph{Say it:} \emph{kadāc ahī̃ bhūkamp āvśe} — Perhaps an earthquake will come here.
  \item \emph{Recall:} say the Gujarati for fire, then the Gujarati for to catch fire, then say \emph{bhūkamp} again
\end{itemize}

\begin{tcolorbox}[breakable,colback=teal!4,colframe=teal!35!black,title={Before you move on}]
What does \gu{ભૂકંપ} mean? (\textquotedblleft{}Earthquake\textquotedblright{}.) Say it once more.
\end{tcolorbox}
